% Original analytic teaching diagram, 2026-11-12.
% A 16-step record of the two-state chain (P(0->1)=0.1, P(1->0)=0.2, started from its invariant law,
% seed 3144): 0000001011111000.  f = indicator of state 1, so hat I = 6/16 = 0.375.
% Batches of length b = 4 give batch means 0, 0.25, 1, 0.25 (red marks); the dashed line is hat I.
\definecolor{VB}{HTML}{7188A3}\definecolor{VBF}{HTML}{E5EAF0}
\definecolor{VG}{HTML}{829C8B}\definecolor{VGF}{HTML}{E8EEE9}
\definecolor{VY}{HTML}{B4A16F}\definecolor{VYF}{HTML}{F1ECDD}
\definecolor{VR}{HTML}{A7787D}\definecolor{VRF}{HTML}{F0E5E6}
\definecolor{VP}{HTML}{9588A5}\definecolor{VPF}{HTML}{ECE8F0}
\definecolor{VI}{HTML}{3B4252}\definecolor{VM}{HTML}{6B7280}\definecolor{VL}{HTML}{4C4D4F}
\begin{tikzpicture}[x=1mm,y=1mm,>=Stealth,font=\figurefont\fontsize{8}{10}\selectfont,
lab/.style={text=VM,font=\figurefont\fontsize{7.5}{9}\selectfont}]
% record: cell t occupies [5.2(t-1), 5.2t] horizontally, gap of 3 mm between batches
\foreach \v [count=\t] in {0,0,0,0,0,0,1,0,1,1,1,1,1,0,0,0}{
  \pgfmathsetmacro{\xx}{5.2*(\t-1)+3*floor((\t-1)/4)}
  \ifnum\v=1
    \filldraw[fill=VG,draw=VG,line width=.6pt] (\xx+.35,20) rectangle (\xx+4.85,24.5);
    \node[text=white] at (\xx+2.6,22.25) {$1$};
  \else
    \filldraw[fill=white,draw=VG,line width=.6pt] (\xx+.35,20) rectangle (\xx+4.85,24.5);
    \node[text=VI] at (\xx+2.6,22.25) {$0$};
  \fi
  \node[lab] at (\xx+2.6,27.3) {$\t$};
}
\node[lab,anchor=east] at (-1,27.3) {$t$};
\node[text=VI,anchor=east] at (-1,22.25) {$f(X^{(t)})$};
% batch braces and means
\foreach \j/\m/\s in {0/0/0,1/0.25/0.25,2/1/1,3/0.25/0.25}{
  \pgfmathsetmacro{\xa}{23.8*\j+.35}
  \pgfmathsetmacro{\xb}{23.8*\j+20.45}
  \draw[VL,line width=.5pt] (\xa,18) -- (\xa,17) -- (\xb,17) -- (\xb,18);
  \node[text=VI] at (23.8*\j+10.4,14) {$\bar Y_{\the\numexpr\j+1\relax}=\s$};
}
% batch means on a scale [0,1] -> [0,10] mm below
\draw[VL,line width=.5pt,->] (-1,0) -- (93,0);
\draw[VL,line width=.5pt] (-1,0) -- (-1,10);
\node[lab,anchor=east] at (-1.5,0) {$0$};
\node[lab,anchor=east] at (-1.5,10) {$1$};
\draw[VL,line width=.4pt] (-1.8,10) -- (-1,10);
\draw[VG,line width=.8pt,dash pattern=on 3pt off 2pt] (0,3.75) -- (91,3.75);
\node[text=VG,anchor=west] at (91.5,3.75) {$\hat I$};
\foreach \j/\m in {0/0,1/0.25,2/1,3/0.25}{
  \draw[VR,line width=1.4pt] (23.8*\j+3,{10*\m}) -- (23.8*\j+17.8,{10*\m});
}
\end{tikzpicture}
